% ch1-espaces-vectoriels.tex — Chapitre 1 : Espaces vectoriels et sous-espaces
% Ce chapitre introduit la structure fondamentale de l'algèbre linéaire :
% l'espace vectoriel. On y développe les notions de sous-espace, de famille
% libre, génératrice, de base, de dimension et de somme directe.

\chapter{Espaces vectoriels et sous-espaces}\label{ch:espaces-vectoriels}

\section*{Introduction}
\addcontentsline{toc}{section}{Introduction}

L'algèbre linéaire est née de l'étude des vecteurs du plan et de
l'espace. Dans $\R^2$, un vecteur $\vec{v}=(x,y)$ peut être
\emph{additionné} à un autre vecteur et \emph{multiplié} par un
scalaire réel, et ces deux opérations satisfont des règles naturelles~:
l'addition est commutative et associative, la multiplication par un
scalaire se distribue sur les sommes, etc. Les mêmes règles s'appliquent
dans $\R^3$.

\medskip

Mais ces opérations ne sont pas l'apanage de la géométrie. Les matrices,
les polynômes, les fonctions continues, les suites réelles --- tous ces
objets peuvent être additionnés entre eux et multipliés par un
scalaire, et les mêmes règles algébriques sont satisfaites. L'idée
fondamentale de l'algèbre linéaire est de \emph{dégager un cadre
abstrait} --- celui d'\emph{espace vectoriel} --- qui unifie tous ces
exemples et permet de les étudier avec les mêmes outils.

\medskip

Ce chapitre constitue le fondement sur lequel repose l'ensemble du
cours. Nous y définissons les espaces vectoriels et leurs sous-espaces,
puis nous développons les notions de famille libre, famille
génératrice, base et dimension. Ces concepts sont indispensables pour
tout ce qui suit~: applications linéaires, matrices, déterminants,
valeurs propres.

\begin{center}
\begin{tikzpicture}[scale=1.2]
  % Axes
  \draw[-{Stealth},gray] (-0.5,0) -- (4,0) node[right] {$x$};
  \draw[-{Stealth},gray] (0,-0.5) -- (0,3.2) node[above] {$y$};
  % Vectors
  \draw[-{Stealth},thick,blue!70!black] (0,0) -- (3,1)
    node[midway,below right] {$\vec{u}$};
  \draw[-{Stealth},thick,red!70!black] (0,0) -- (1,2.5)
    node[midway,above left] {$\vec{v}$};
  \draw[-{Stealth},thick,purple!70!black,dashed] (0,0) -- (2.5,2)
    node[midway,above,sloped,font=\small] {$\alpha\vec{u}+\beta\vec{v}$};
  % Parallelogram
  \draw[dotted,gray] (3,1) -- (4,3.5);
  \draw[dotted,gray] (1,2.5) -- (4,3.5);
  % Sum vector
  \draw[-{Stealth},thick,green!50!black] (0,0) -- (4,3.5)
    node[right] {$\vec{u}+\vec{v}$};
\end{tikzpicture}
\captionof{figure}{Combinaison linéaire de deux vecteurs dans $\R^2$.}
\label{fig:comb-lin-R2}
\end{center}

% ============================================================
\section{Espace vectoriel sur un corps}\label{sec:ev-definition}
% ============================================================

\subsection{Définition axiomatique}\label{subsec:ev-axiomes}

Dans tout ce chapitre, $\K$ désigne un corps (voir 
ef{subsec:corps}).
Les éléments de $\K$ sont appelés \emph{scalaires}\index{scalaire}. En
pratique, le lecteur pourra penser à $\K=\R$ ou $\K=\C$.

\begin{definition}{Espace vectoriel}{def:ev}
  Un \emph{espace vectoriel}\index{espace vectoriel} sur $\K$ (ou
  $\K$-espace vectoriel) est un triplet $(E,+,\cdot)$ où $E$ est un
  ensemble non vide, $+\colon E\times E\to E$ est une loi de composition
  interne (l'\emph{addition vectorielle}) et
  $\cdot\colon\K\times E\to E$ est une loi de composition externe (la
  \emph{multiplication par un scalaire}\index{multiplication scalaire}),
  satisfaisant les axiomes suivants~:

  \medskip
  \noindent\textbf{Axiomes d'addition.}
  \begin{enumerate}[label=\textbf{(EV\arabic*)}]
    \item\label{ax:ev-assoc} \textbf{Associativité~:}
      $\forall\, u,v,w\in E,\; (u+v)+w = u+(v+w)$.
    \item\label{ax:ev-comm} \textbf{Commutativité~:}
      $\forall\, u,v\in E,\; u+v = v+u$.
    \item\label{ax:ev-neutre} \textbf{Élément neutre~:}
      $\exists\, \vzero\in E,\;\forall\, u\in E,\; u+\vzero = u$.
    \item\label{ax:ev-oppose} \textbf{Opposé~:}
      $\forall\, u\in E,\;\exists\, (-u)\in E,\; u+(-u) = \vzero$.
  \end{enumerate}
  Autrement dit, $(E,+)$ est un groupe abélien.

  \medskip
  \noindent\textbf{Axiomes de multiplication scalaire.}
  \begin{enumerate}[label=\textbf{(EV\arabic*)}]
    \setcounter{enumi}{4}
    \item\label{ax:ev-compat} \textbf{Compatibilité~:}
      $\forall\, \alpha,\beta\in\K,\;\forall\, u\in E,\;
      \alpha\cdot(\beta\cdot u) = (\alpha\beta)\cdot u$.
    \item\label{ax:ev-unite} \textbf{Élément unité~:}
      $\forall\, u\in E,\; 1_\K\cdot u = u$.
    \item\label{ax:ev-dist1} \textbf{Distributivité (scalaire)~:}
      $\forall\, \alpha\in\K,\;\forall\, u,v\in E,\;
      \alpha\cdot(u+v) = \alpha\cdot u + \alpha\cdot v$.
    \item\label{ax:ev-dist2} \textbf{Distributivité (vecteur)~:}
      $\forall\, \alpha,\beta\in\K,\;\forall\, u\in E,\;
      (\alpha+\beta)\cdot u = \alpha\cdot u + \beta\cdot u$.
  \end{enumerate}
  Les éléments de $E$ sont appelés \emph{vecteurs}\index{vecteur}.
\end{definition}

\begin{remark}{Notation}{rem:notation-mult-scal}
  On omet généralement le point $\cdot$ et on écrit $\alpha u$ au lieu
  de $\alpha\cdot u$. La multiplication scalaire n'est \emph{pas} une
  loi de composition interne~: elle prend un scalaire dans $\K$ et un
  vecteur dans $E$, et renvoie un vecteur dans $E$.
\end{remark}

\begin{proposition}{Propriétés immédiates}{prop:ev-proprietes}
  Soit $E$ un $\K$-espace vectoriel. Pour tous $\alpha\in\K$ et $u\in E$~:
  \begin{enumerate}[label=(\roman*)]
    \item $0_\K\cdot u = \vzero$.
    \item $\alpha\cdot\vzero = \vzero$.
    \item $(-1_\K)\cdot u = -u$.
    \item $\alpha\cdot u = \vzero \implies \alpha = 0_\K \text{ ou } u = \vzero$.
  \end{enumerate}
\end{proposition}

\begin{proof}
  \emph{(i)} On a
  $0_\K\cdot u = (0_\K + 0_\K)\cdot u = 0_\K\cdot u + 0_\K\cdot u$
  par \ref{ax:ev-dist2}. En ajoutant l'opposé de $0_\K\cdot u$ des
  deux côtés, on obtient $\vzero = 0_\K\cdot u$.

  \emph{(ii)} On a
  $\alpha\cdot\vzero = \alpha\cdot(\vzero+\vzero) =
  \alpha\cdot\vzero + \alpha\cdot\vzero$ par \ref{ax:ev-dist1}, d'où
  $\vzero = \alpha\cdot\vzero$ en simplifiant.

  \emph{(iii)} On calcule
  $u + (-1_\K)\cdot u = 1_\K\cdot u + (-1_\K)\cdot u =
  (1_\K + (-1_\K))\cdot u = 0_\K\cdot u = \vzero$ par (i). Donc
  $(-1_\K)\cdot u$ est l'opposé de $u$.

  \emph{(iv)} Supposons $\alpha\neq 0_\K$. Alors $\alpha^{-1}$ existe
  dans $\K$, et
  $u = 1_\K\cdot u = (\alpha^{-1}\alpha)\cdot u =
  \alpha^{-1}\cdot(\alpha\cdot u) = \alpha^{-1}\cdot\vzero = \vzero$
  par (ii).
\end{proof}

% ============================================================
\subsection{Exemples fondamentaux}\label{subsec:ev-exemples}
% ============================================================

\begin{example}{L'espace $\K^n$}{ex:Kn}
  \index{$\K^n$}%
  L'ensemble des $n$-uplets
  \[
    \K^n = \setcond{(x_1,\ldots,x_n)}{x_i\in\K}
  \]
  muni de l'addition composante par composante
  \[
    (x_1,\ldots,x_n)+(y_1,\ldots,y_n) = (x_1+y_1,\ldots,x_n+y_n)
  \]
  et de la multiplication scalaire
  \[
    \alpha(x_1,\ldots,x_n) = (\alpha x_1,\ldots,\alpha x_n)
  \]
  est un $\K$-espace vectoriel. Le vecteur nul est
  $\vzero=(0,\ldots,0)$. C'est l'exemple le plus fondamental~; pour
  $n=2$ et $n=3$, il correspond aux vecteurs du plan et de l'espace.
\end{example}

\begin{example}{L'espace des matrices $\Mat_{n,p}(\K)$}{ex:matrices}
  \index{espace des matrices}%
  L'ensemble $\Mat_{n,p}(\K)$ des matrices à $n$ lignes et $p$ colonnes
  à coefficients dans $\K$, muni de l'addition matricielle et de la
  multiplication par un scalaire, est un $\K$-espace vectoriel. Le
  vecteur nul est la matrice nulle $0_{n\times p}$.
\end{example}

\begin{example}{L'espace des polynômes}{ex:polynomes}
  \index{espace des polynômes}%
  \begin{enumerate}[label=(\alph*)]
    \item L'ensemble $\K[X]$ de tous les polynômes à coefficients dans $\K$,
      muni de l'addition des polynômes et de la multiplication par un
      scalaire, est un $\K$-espace vectoriel.
    \item Pour tout $n\in\N$, le sous-ensemble
      \[
        \K_n[X] \deff \setcond{P\in\K[X]}{\deg P\leq n}\cup\set{0}
      \]
      des polynômes de degré au plus $n$ (incluant le polynôme nul) est
      un $\K$-espace vectoriel.
  \end{enumerate}
\end{example}

\begin{example}{Espaces de fonctions}{ex:fonctions}
  \index{espace de fonctions}%
  Soit $X$ un ensemble non vide. L'ensemble
  $\cF(X,\K) \deff \setcond{f}{f\colon X\to\K}$ de toutes les
  fonctions de $X$ dans $\K$, muni des opérations
  \[
    (f+g)(x)=f(x)+g(x), \qquad (\alpha f)(x) = \alpha f(x),
  \]
  est un $\K$-espace vectoriel. De même, si $\K=\R$, les ensembles
  suivants sont des $\R$-espaces vectoriels~:
  \begin{itemize}
    \item $\cC^0(\R,\R)$, l'espace des fonctions continues~;
    \item $\cC^k(\R,\R)$, l'espace des fonctions $k$ fois dérivables
      à dérivée $k$-ième continue~;
    \item $\cC^\infty(\R,\R)$, l'espace des fonctions indéfiniment
      dérivables.
  \end{itemize}
\end{example}

\begin{example}{Espace des solutions d'un système homogène}{ex:solutions-homogene}
  \index{système homogène}%
  Soit $A\in\Mat_{n,p}(\K)$. L'ensemble des solutions du système
  linéaire homogène $AX=0$~:
  \[
    \cS = \setcond{X\in\K^p}{AX=0}
  \]
  est un $\K$-espace vectoriel (c'est un sous-espace de $\K^p$,
  comme nous le verrons dans la 
ef{sec:sous-espaces}).
\end{example}

\begin{example}{L'espace des suites}{ex:suites}
  \index{espace des suites}%
  L'ensemble $\K^\N$ de toutes les suites $(u_n)_{n\geq 0}$ à valeurs
  dans $\K$, muni de l'addition terme à terme et de la multiplication
  scalaire, est un $\K$-espace vectoriel.
\end{example}

\begin{remark}{Contre-exemples}{rem:non-ev}
  L'ensemble $\N$ muni de l'addition usuelle n'est \emph{pas} un espace
  vectoriel sur $\R$~: d'une part, il n'y a pas d'opposé ($-3\notin\N$),
  et d'autre part, la multiplication scalaire par un réel ne reste pas
  dans $\N$ ($\frac{1}{2}\cdot 3 = \frac{3}{2}\notin\N$).
\end{remark}

% ============================================================
\section{Sous-espaces vectoriels}\label{sec:sous-espaces}
% ============================================================

\begin{definition}{Sous-espace vectoriel}{def:sev}
  Soit $E$ un $\K$-espace vectoriel. Une partie $F\subseteq E$ est un
  \emph{sous-espace vectoriel}\index{sous-espace vectoriel} de $E$
  (en abrégé~: \emph{s.e.v.}) si $F$, muni des opérations induites par
  $E$, est lui-même un $\K$-espace vectoriel.
\end{definition}

La vérification directe de tous les axiomes est lourde. Le théorème
suivant fournit un critère simple et efficace.

\begin{theorem}{Caractérisation des sous-espaces vectoriels}{thm:caract-sev}
  Soit $E$ un $\K$-espace vectoriel et $F\subseteq E$. Alors $F$ est
  un sous-espace vectoriel de $E$ si et seulement si~:
  \begin{enumerate}[label=(\roman*)]
    \item $F\neq\emptyset$ (ou, de manière équivalente, $\vzero\in F$)~;
    \item $\forall\, u,v\in F,\; u+v\in F$ \emph{(stabilité par addition)}~;
    \item $\forall\, \alpha\in\K,\;\forall\, u\in F,\; \alpha u\in F$
      \emph{(stabilité par multiplication scalaire)}.
  \end{enumerate}
  De manière équivalente, les conditions (ii) et (iii) peuvent être
  remplacées par la condition unique~:
  \begin{enumerate}[label=(\roman*')]
    \setcounter{enumi}{1}
    \item $\forall\, \alpha,\beta\in\K,\;\forall\, u,v\in F,\;
      \alpha u + \beta v\in F$
      \emph{(stabilité par combinaison linéaire)}.
  \end{enumerate}
\end{theorem}

\begin{proof}
  $(\Rightarrow)$ Si $F$ est un s.e.v.\ de $E$, alors $F$ contient son
  vecteur nul, qui doit être $\vzero$ (car $\vzero_F = 0_\K\cdot u =
  \vzero_E$ pour tout $u\in F$). Les propriétés (ii) et (iii) découlent
  du fait que $+$ et $\cdot$ sont des lois sur $F$.

  $(\Leftarrow)$ Supposons (i), (ii), (iii). Puisque $F\neq\emptyset$,
  il existe $u\in F$. Par (iii), $0_\K\cdot u = \vzero\in F$. Par
  (iii), $(-1_\K)\cdot u = -u\in F$ pour tout $u\in F$, donc tout
  vecteur de $F$ admet un opposé dans $F$. Par (ii), $F$ est stable
  par addition. Les axiomes \ref{ax:ev-assoc}--\ref{ax:ev-dist2} sont
  satisfaits dans $F$ car ils le sont dans $E$ et $F\subseteq E$.
  Ainsi $(F,+,\cdot)$ est un $\K$-espace vectoriel.

  L'équivalence entre (ii)+(iii) et (ii') est immédiate~: (ii') avec
  $\alpha=\beta=1$ donne (ii), et (ii') avec $\beta=0$ donne (iii)~;
  réciproquement, si (ii) et (iii) sont vérifiées, alors $\alpha u\in F$
  et $\beta v\in F$ par (iii), puis $\alpha u+\beta v\in F$ par (ii).
\end{proof}

\begin{example}{Sous-espaces de $\R^3$}{ex:sev-R3}
  Les sous-espaces vectoriels de $\R^3$ sont~:
  \begin{itemize}
    \item $\set{\vzero}$, le sous-espace nul (un point)~;
    \item les droites vectorielles passant par l'origine~;
    \item les plans vectoriels passant par l'origine~;
    \item $\R^3$ lui-même.
  \end{itemize}
  Par exemple, $D=\setcond{t(1,2,3)}{t\in\R}$ est une droite
  vectorielle, et $P=\setcond{(x,y,z)\in\R^3}{x+y+z=0}$ est un plan
  vectoriel.
\end{example}

\begin{center}
\begin{tikzpicture}[scale=0.9,
  x={({cos(210)*0.8cm},{sin(210)*0.8cm})},
  y={(1cm,0cm)},
  z={(0cm,1cm)}]
  % Axes
  \draw[-{Stealth},gray] (0,0,0) -- (3,0,0) node[below left] {$x$};
  \draw[-{Stealth},gray] (0,0,0) -- (0,3.5,0) node[right] {$y$};
  \draw[-{Stealth},gray] (0,0,0) -- (0,0,3) node[above] {$z$};
  % Plane x+y+z=0
  \fill[blue!15,opacity=0.5]
    (2,-1,-1) -- (-1,2,-1) -- (-1,-1,2) -- cycle;
  \draw[blue!50,thick]
    (2,-1,-1) -- (-1,2,-1) -- (-1,-1,2) -- cycle;
  \node[blue!70!black,font=\small] at (0.5,1.2,-0.8) {$x{+}y{+}z=0$};
  % Line through (1,1,1)
  \draw[red!70!black,thick,-{Stealth}] (0,0,0) -- (1.5,1.5,1.5);
  \draw[red!70!black,thick,dashed] (0,0,0) -- (-0.8,-0.8,-0.8);
  \node[red!70!black,font=\small,right] at (1.5,1.5,1.5) {$\Vect(1,1,1)$};
  % Origin
  \fill (0,0,0) circle (1.5pt);
\end{tikzpicture}
\captionof{figure}{Un plan vectoriel et une droite vectorielle dans $\R^3$.}
\label{fig:sev-R3}
\end{center}

\begin{example}{Sous-espaces de $\K[X]$}{ex:sev-polynomes}
  Soit $n\in\N$. L'ensemble $\K_n[X]$ des polynômes de degré $\leq n$
  est un sous-espace vectoriel de $\K[X]$~:
  \begin{itemize}
    \item $\vzero = 0\in\K_n[X]$~;
    \item si $\deg P\leq n$ et $\deg Q\leq n$, alors
      $\deg(P+Q)\leq n$~;
    \item si $\deg P\leq n$, alors $\deg(\alpha P)\leq n$ pour tout
      $\alpha\in\K$.
  \end{itemize}
  En revanche, $\setcond{P\in\K[X]}{\deg P = n}$ n'est \emph{pas} un
  sous-espace~: il ne contient pas le polynôme nul, et n'est pas stable
  par addition (la somme de deux polynômes de degré $n$ peut avoir un
  degré inférieur).
\end{example}

\begin{example}{Noyau d'une application linéaire}{ex:sev-noyau}
  Nous verrons au 
ef{ch:espaces-vectoriels} que le noyau et l'image
  d'une application linéaire sont des sous-espaces vectoriels. Pour
  l'instant, vérifions-le dans un cas particulier~: soit
  $\vphi\colon\R^3\to\R$ définie par $\vphi(x,y,z)=2x-y+3z$. Alors
  \[
    \Ker\vphi = \setcond{(x,y,z)\in\R^3}{2x-y+3z=0}
  \]
  est un sous-espace vectoriel de $\R^3$ (c'est le plan
  d'équation $2x-y+3z=0$). En effet, $\vzero\in\Ker\vphi$, et si
  $u,v\in\Ker\vphi$ et $\alpha\in\R$, alors
  $\vphi(\alpha u+v)=\alpha\vphi(u)+\vphi(v)=0$, donc
  $\alpha u+v\in\Ker\vphi$.
\end{example}

\begin{proposition}{Propriétés élémentaires des sous-espaces}{prop:sev-elem}
  Soit $E$ un $\K$-espace vectoriel.
  \begin{enumerate}[label=(\roman*)]
    \item $\set{\vzero}$ et $E$ sont des sous-espaces vectoriels de $E$
      (appelés \emph{sous-espaces triviaux}\index{sous-espace!trivial}).
    \item Tout sous-espace vectoriel $F$ de $E$ contient $\vzero$.
    \item Si $F$ est un s.e.v.\ de $E$ et $G$ est un s.e.v.\ de $F$,
      alors $G$ est un s.e.v.\ de $E$.
  \end{enumerate}
\end{proposition}

\begin{proof}
  Les points (i) et (ii) sont immédiats d'après le
  
ef{thm:thm:caract-sev}. Pour (iii), $G$ est non vide et stable par
  combinaison linéaire dans $F$, donc aussi dans $E$ puisque les
  opérations sont les mêmes.
\end{proof}

% ============================================================
\section{Intersection et somme de sous-espaces}\label{sec:intersection-somme}
% ============================================================

\subsection{Intersection}\label{subsec:intersection}

\begin{proposition}{Intersection de sous-espaces}{prop:intersection-sev}
  \index{intersection de sous-espaces}%
  Soit $E$ un $\K$-espace vectoriel et $(F_i)_{i\in I}$ une famille de
  sous-espaces vectoriels de $E$. Alors l'intersection
  $\displaystyle\bigcap_{i\in I}F_i$ est un sous-espace vectoriel de $E$.
\end{proposition}

\begin{proof}
  Posons $F=\bigcap_{i\in I}F_i$. Puisque $\vzero\in F_i$ pour tout
  $i\in I$, on a $\vzero\in F$, donc $F\neq\emptyset$. Soient
  $u,v\in F$ et $\alpha,\beta\in\K$. Pour tout $i\in I$, $u,v\in F_i$,
  donc $\alpha u+\beta v\in F_i$ (car $F_i$ est un s.e.v.). Ainsi
  $\alpha u+\beta v\in F$. Par le 
ef{thm:thm:caract-sev}, $F$ est un
  s.e.v.\ de $E$.
\end{proof}

\begin{remark}{La réunion n'est pas un sous-espace}{rem:reunion-sev}
  En général, la réunion de deux sous-espaces vectoriels n'est
  \emph{pas} un sous-espace vectoriel. Par exemple, dans $\R^2$, soient
  $F=\setcond{(t,0)}{t\in\R}$ (l'axe des abscisses) et
  $G=\setcond{(0,t)}{t\in\R}$ (l'axe des ordonnées). Alors
  $(1,0)\in F\cup G$ et $(0,1)\in F\cup G$, mais
  $(1,0)+(0,1)=(1,1)\notin F\cup G$.
\end{remark}

\subsection{Somme de sous-espaces}\label{subsec:somme}

\begin{definition}{Somme de sous-espaces}{def:somme-sev}
  \index{somme de sous-espaces}%
  Soient $F$ et $G$ deux sous-espaces vectoriels d'un $\K$-espace
  vectoriel $E$. La \emph{somme} de $F$ et $G$ est l'ensemble
  \[
    F+G \deff \setcond{u+v}{u\in F,\;v\in G}.
  \]
  Plus généralement, si $F_1,\ldots,F_r$ sont des sous-espaces de $E$,
  \[
    F_1+\cdots+F_r = \setcond{u_1+\cdots+u_r}{u_i\in F_i}.
  \]
\end{definition}

\begin{proposition}{La somme est un sous-espace}{prop:somme-sev}
  $F+G$ est un sous-espace vectoriel de $E$. De plus, c'est le plus
  petit sous-espace de $E$ contenant à la fois $F$ et $G$.
\end{proposition}

\begin{proof}
  \emph{$F+G$ est un s.e.v.} Le vecteur $\vzero=\vzero+\vzero\in F+G$,
  donc $F+G\neq\emptyset$. Soient $w_1=u_1+v_1$ et $w_2=u_2+v_2$
  dans $F+G$ (avec $u_i\in F$, $v_i\in G$) et $\alpha,\beta\in\K$.
  Alors
  \[
    \alpha w_1+\beta w_2 =
    \underbrace{(\alpha u_1+\beta u_2)}_{\in F}+
    \underbrace{(\alpha v_1+\beta v_2)}_{\in G}\in F+G.
  \]

  \emph{Minimalité.} On a $F\subseteq F+G$ (car tout $u\in F$ s'écrit
  $u=u+\vzero$) et $G\subseteq F+G$. Si $H$ est un s.e.v.\ contenant
  $F$ et $G$, alors pour tout $u\in F$ et $v\in G$, on a $u,v\in H$,
  donc $u+v\in H$, d'où $F+G\subseteq H$.
\end{proof}

% ============================================================
\section{Somme directe et sous-espaces supplémentaires}\label{sec:somme-directe}
% ============================================================

\begin{definition}{Somme directe (interne)}{def:somme-directe}
  \index{somme directe}%
  Soient $F$ et $G$ deux sous-espaces vectoriels de $E$. On dit que la
  somme $F+G$ est \emph{directe} si tout vecteur $w\in F+G$ s'écrit de
  manière \emph{unique} sous la forme $w=u+v$ avec $u\in F$ et $v\in G$.
  On note alors $F\oplus G$\index{$\oplus$}.
\end{definition}

\begin{theorem}{Caractérisation de la somme directe}{thm:caract-somme-directe}
  Soient $F$ et $G$ deux sous-espaces vectoriels de $E$. Les assertions
  suivantes sont équivalentes~:
  \begin{enumerate}[label=(\roman*)]
    \item La somme $F+G$ est directe.
    \item $F\cap G = \set{\vzero}$.
    \item Pour tous $u\in F$ et $v\in G$, si $u+v=\vzero$, alors
      $u=\vzero$ et $v=\vzero$.
  \end{enumerate}
\end{theorem}

\begin{proof}
  \emph{(i)$\Rightarrow$(ii).} Soit $w\in F\cap G$. Alors $w$
  s'écrit $w=w+\vzero$ (avec $w\in F$, $\vzero\in G$) et
  $w=\vzero+w$ (avec $\vzero\in F$, $w\in G$). Par unicité de la
  décomposition, $w=\vzero$.

  \emph{(ii)$\Rightarrow$(iii).} Si $u+v=\vzero$ avec $u\in F$,
  $v\in G$, alors $u=-v\in G$. Donc $u\in F\cap G=\set{\vzero}$, d'où
  $u=\vzero$ et $v=\vzero$.

  \emph{(iii)$\Rightarrow$(i).} Supposons $w=u_1+v_1=u_2+v_2$ avec
  $u_i\in F$, $v_i\in G$. Alors
  $(u_1-u_2)+(v_1-v_2)=\vzero$ avec $u_1-u_2\in F$ et
  $v_1-v_2\in G$. Par hypothèse, $u_1-u_2=\vzero$ et
  $v_1-v_2=\vzero$, donc $u_1=u_2$ et $v_1=v_2$. La décomposition est
  unique.
\end{proof}

\begin{definition}{Sous-espaces supplémentaires}{def:supplementaires}
  \index{sous-espaces supplémentaires}%
  Deux sous-espaces $F$ et $G$ de $E$ sont dits \emph{supplémentaires}
  (dans $E$) si $E=F\oplus G$, c'est-à-dire si~:
  \[
    E = F+G \qquad\text{et}\qquad F\cap G = \set{\vzero}.
  \]
  On dit aussi que $G$ est un \emph{supplémentaire}\index{supplémentaire}
  de $F$ dans $E$.
\end{definition}

\begin{example}{Supplémentaires dans $\R^3$}{ex:supplementaires-R3}
  Soient
  $F=\setcond{(x,y,0)}{x,y\in\R}$ (le plan $z=0$) et
  $G=\setcond{(0,0,z)}{z\in\R}$ (l'axe $Oz$). Alors
  $\R^3=F\oplus G$~: tout $(x,y,z)\in\R^3$ s'écrit de manière unique
  comme $(x,y,0)+(0,0,z)$, et $F\cap G=\set{(0,0,0)}$.
\end{example}

\begin{remark}{Non-unicité du supplémentaire}{rem:supplementaire-non-unique}
  Un sous-espace admet en général \emph{plusieurs} supplémentaires.
  Par exemple, dans $\R^2$, la droite $F=\Vect((1,0))$ admet pour
  supplémentaire toute droite $G=\Vect((a,b))$ avec $b\neq 0$.
\end{remark}

La notion de somme directe se généralise à un nombre fini quelconque de
sous-espaces.

\begin{definition}{Somme directe de plusieurs sous-espaces}{def:somme-directe-n}
  Soient $F_1,\ldots,F_r$ des sous-espaces de $E$. On dit que la somme
  $F_1+\cdots+F_r$ est \emph{directe}, et on note
  $F_1\oplus\cdots\oplus F_r$, si tout vecteur $w\in F_1+\cdots+F_r$
  s'écrit de manière unique $w=u_1+\cdots+u_r$ avec $u_i\in F_i$.
\end{definition}

\begin{proposition}{Caractérisation de la somme directe multiple}{prop:caract-somme-directe-n}
  La somme $F_1+\cdots+F_r$ est directe si et seulement si pour tout
  $j\in\set{1,\ldots,r}$~:
  \[
    F_j \cap \bigl(F_1+\cdots+F_{j-1}+F_{j+1}+\cdots+F_r\bigr) =
    \set{\vzero}.
  \]
  De manière équivalente, si $u_1+\cdots+u_r=\vzero$ avec $u_i\in F_i$,
  alors $u_1=\cdots=u_r=\vzero$.
\end{proposition}

\begin{proof}
  La preuve suit le même schéma que celle du 
ef{thm:thm:caract-somme-directe}.
  Montrons l'équivalence avec la condition sur l'unicité.

  $(\Rightarrow)$ Si $u_1+\cdots+u_r=\vzero$ avec $u_i\in F_i$, on a
  aussi $\vzero+\cdots+\vzero=\vzero$. Par unicité de la
  décomposition, $u_i=\vzero$ pour tout $i$.

  $(\Leftarrow)$ Si $u_1+\cdots+u_r=u_1'+\cdots+u_r'$ avec
  $u_i,u_i'\in F_i$, alors $(u_1-u_1')+\cdots+(u_r-u_r')=\vzero$
  avec $u_i-u_i'\in F_i$. Par hypothèse, $u_i-u_i'=\vzero$ pour tout
  $i$, soit $u_i=u_i'$.
\end{proof}

% ============================================================
\section{Combinaisons linéaires et sous-espace engendré}\label{sec:comb-lin}
% ============================================================

\begin{definition}{Combinaison linéaire}{def:comb-lin}
  \index{combinaison linéaire}%
  Soit $E$ un $\K$-espace vectoriel. Soient $v_1,\ldots,v_n\in E$ et
  $\alpha_1,\ldots,\alpha_n\in\K$. Le vecteur
  \[
    w = \alpha_1 v_1 + \alpha_2 v_2 + \cdots + \alpha_n v_n
    = \sum_{i=1}^{n}\alpha_i v_i
  \]
  est appelé \emph{combinaison linéaire} des vecteurs $v_1,\ldots,v_n$
  à coefficients $\alpha_1,\ldots,\alpha_n$.
\end{definition}

\begin{definition}{Sous-espace engendré}{def:sous-espace-engendre}
  \index{sous-espace engendré}\index{Vect@$\Vect$}%
  Soit $S\subseteq E$ une partie (éventuellement infinie). Le
  \emph{sous-espace engendré par $S$}, noté $\Vect(S)$, est le plus
  petit sous-espace vectoriel de $E$ contenant $S$. De manière
  équivalente~:
  \[
    \Vect(S) = \setcond{\alpha_1 v_1+\cdots+\alpha_n v_n}{%
      n\in\N,\; v_i\in S,\; \alpha_i\in\K}.
  \]
  Par convention, $\Vect(\emptyset)=\set{\vzero}$.

  Si $S=\set{v_1,\ldots,v_n}$ est finie, on note
  $\Vect(v_1,\ldots,v_n)$.
\end{definition}

\begin{proposition}{$\Vect(S)$ est l'intersection de tous les
  sous-espaces contenant $S$}{prop:vect-intersection}
  Pour toute partie $S\subseteq E$~:
  \[
    \Vect(S) = \bigcap_{\substack{F\text{ s.e.v.\ de }E \\ S\subseteq F}} F.
  \]
\end{proposition}

\begin{proof}
  Notons $V$ l'intersection du membre droit. Par la
  
ef{prop:prop:intersection-sev}, $V$ est un s.e.v.\ de $E$, et
  $S\subseteq V$ (car $S\subseteq F$ pour chaque $F$ dans
  l'intersection). Donc $V$ est un s.e.v.\ contenant $S$. Par
  définition, $\Vect(S)\subseteq V$ (car $\Vect(S)$ est le plus petit
  tel s.e.v.).

  Inversement, $\Vect(S)$ est un s.e.v.\ contenant $S$, donc il
  apparaît parmi les $F$ de l'intersection, d'où
  $V\subseteq\Vect(S)$.
\end{proof}

\begin{example}{Sous-espace engendré dans $\R^3$}{ex:vect-R3}
  Soient $v_1=(1,0,1)$ et $v_2=(0,1,-1)$. Alors
  \[
    \Vect(v_1,v_2) = \setcond{\alpha(1,0,1)+\beta(0,1,-1)}{\alpha,\beta\in\R}
    = \setcond{(\alpha,\beta,\alpha-\beta)}{\alpha,\beta\in\R}.
  \]
  C'est un plan passant par l'origine dans $\R^3$, d'équation
  $x-y-z=0$ (on vérifie que $\alpha-\beta-(\alpha-\beta)=0$).
\end{example}

\begin{definition}{Famille génératrice}{def:famille-generatrice}
  \index{famille génératrice}%
  Soit $E$ un $\K$-espace vectoriel. Une famille $(v_i)_{i\in I}$ de
  vecteurs de $E$ est dite \emph{génératrice} de $E$ si
  $\Vect((v_i)_{i\in I})=E$, c'est-à-dire si tout vecteur de $E$ est
  combinaison linéaire de la famille.
\end{definition}

\begin{example}{Famille génératrice de $\R^3$}{ex:generatrice-R3}
  La famille $((1,0,0),(0,1,0),(0,0,1))$ est génératrice de $\R^3$~:
  tout $(x,y,z)\in\R^3$ s'écrit $x(1,0,0)+y(0,1,0)+z(0,0,1)$. La
  famille $((1,1,0),(1,0,1),(0,1,1))$ est également génératrice
  (le vérifier en exercice).
\end{example}

% ============================================================
\section{Familles libres}\label{sec:familles-libres}
% ============================================================

\begin{definition}{Famille libre (linéairement indépendante)}{def:famille-libre}
  \index{famille libre}\index{indépendance linéaire}%
  Une famille $(v_1,\ldots,v_n)$ de vecteurs de $E$ est dite
  \emph{libre} (ou \emph{linéairement indépendante}) si la seule
  combinaison linéaire nulle est la combinaison triviale~:
  \[
    \alpha_1 v_1+\cdots+\alpha_n v_n = \vzero
    \implies
    \alpha_1=\cdots=\alpha_n=0_\K.
  \]
  Une famille qui n'est pas libre est dite \emph{liée}%
  \index{famille liée} (ou \emph{linéairement dépendante}).
\end{definition}

\begin{remark}{Famille liée}{rem:famille-liee}
  La famille $(v_1,\ldots,v_n)$ est liée si et seulement s'il existe
  des scalaires $\alpha_1,\ldots,\alpha_n\in\K$, \emph{non tous nuls},
  tels que $\alpha_1 v_1+\cdots+\alpha_n v_n=\vzero$. De manière
  équivalente, l'un des vecteurs est combinaison linéaire des autres.
\end{remark}

\begin{proposition}{Propriétés des familles libres}{prop:libre-proprietes}
  \begin{enumerate}[label=(\roman*)]
    \item La famille vide est libre (par convention).
    \item Toute sous-famille d'une famille libre est libre.
    \item Toute famille contenant le vecteur nul $\vzero$ est liée.
    \item Toute famille contenant un vecteur en double est liée.
    \item Une famille $(v)$ réduite à un seul vecteur est libre si et
      seulement si $v\neq\vzero$.
    \item Une famille $(v_1,v_2)$ est liée si et seulement si $v_1$ et
      $v_2$ sont colinéaires (l'un est multiple scalaire de l'autre).
  \end{enumerate}
\end{proposition}

\begin{proof}
  \emph{(i)} Par convention, il n'existe pas de combinaison linéaire
  non triviale de la famille vide.

  \emph{(ii)} Si une sous-famille admet une relation de dépendance
  linéaire non triviale, on peut la compléter en ajoutant des
  coefficients $0$ pour les vecteurs absents, ce qui donne une relation
  de dépendance non triviale de la famille entière, contredisant sa
  liberté.

  \emph{(iii)} Si $v_j=\vzero$, la relation $0\cdot v_1+\cdots+1\cdot
  v_j+\cdots+0\cdot v_n=\vzero$ est non triviale.

  \emph{(iv)} Si $v_i=v_j$ avec $i\neq j$, la relation avec
  $\alpha_i=1$, $\alpha_j=-1$ et les autres coefficients nuls est non
  triviale.

  \emph{(v)} $\alpha v=\vzero$ implique $\alpha=0$ ou $v=\vzero$. Donc
  $(v)$ est libre si et seulement si $v\neq\vzero$.

  \emph{(vi)} $(v_1,v_2)$ est liée ssi il existe $(\alpha_1,\alpha_2)
  \neq(0,0)$ avec $\alpha_1 v_1+\alpha_2 v_2=\vzero$. Si $\alpha_1
  \neq 0$, alors $v_1=-(\alpha_2/\alpha_1)v_2$. Si $\alpha_2\neq 0$,
  alors $v_2=-(\alpha_1/\alpha_2)v_1$.
\end{proof}

\begin{example}{Famille libre dans $\R^3$}{ex:libre-R3}
  Montrons que $\cF=(e_1,e_2,e_3)$ avec $e_1=(1,0,0)$, $e_2=(0,1,0)$,
  $e_3=(0,0,1)$ est libre. Supposons
  \[
    \alpha_1(1,0,0)+\alpha_2(0,1,0)+\alpha_3(0,0,1)=(0,0,0).
  \]
  Alors $(\alpha_1,\alpha_2,\alpha_3)=(0,0,0)$, d'où la liberté.
\end{example}

\begin{example}{Famille libre de polynômes}{ex:libre-polynomes}
  Dans $\R[X]$, la famille $(1,X,X^2,\ldots,X^n)$ est libre~: si
  $\alpha_0+\alpha_1 X+\cdots+\alpha_n X^n=0$ (le polynôme nul),
  alors $\alpha_0=\alpha_1=\cdots=\alpha_n=0$ par identification des
  coefficients.
\end{example}

\begin{theorem}{Caractérisation des familles liées}{thm:caract-liee}
  La famille $(v_1,\ldots,v_n)$ ($n\geq 2$) est liée si et seulement
  si l'un des vecteurs est combinaison linéaire des \emph{précédents}~:
  il existe $j\in\set{2,\ldots,n}$ tel que
  $v_j\in\Vect(v_1,\ldots,v_{j-1})$.
\end{theorem}

\begin{proof}
  $(\Leftarrow)$ Immédiat~: si $v_j=\sum_{i=1}^{j-1}\alpha_i v_i$,
  alors $\sum_{i=1}^{j-1}\alpha_i v_i - v_j = \vzero$ est une
  relation non triviale.

  $(\Rightarrow)$ Soit $\alpha_1 v_1+\cdots+\alpha_n v_n=\vzero$ une
  relation non triviale. Soit $j$ le plus grand indice tel que
  $\alpha_j\neq 0$. Si $j=1$, alors $\alpha_1 v_1=\vzero$ avec
  $\alpha_1\neq 0$, donc $v_1=\vzero\in\Vect(\emptyset)$, ce qui
  n'est pas tout à fait la forme souhaitée. Cependant, si $v_1=\vzero$,
  la famille est liée de manière triviale. Pour $j\geq 2$, on a
  $\alpha_j v_j = -\sum_{i=1}^{j-1}\alpha_i v_i$, d'où
  $v_j = -\sum_{i=1}^{j-1}(\alpha_i/\alpha_j)v_i\in
  \Vect(v_1,\ldots,v_{j-1})$.
\end{proof}

% ============================================================
\section{Bases}\label{sec:bases}
% ============================================================

\begin{definition}{Base}{def:base}
  \index{base}%
  Une famille $(v_1,\ldots,v_n)$ de vecteurs d'un $\K$-espace vectoriel
  $E$ est une \emph{base} de $E$ si elle est à la fois
  \emph{libre} et \emph{génératrice}.
\end{definition}

\begin{theorem}{Caractérisation d'une base}{thm:caract-base}
  Soit $\cB=(v_1,\ldots,v_n)$ une famille de vecteurs de $E$. Les
  assertions suivantes sont équivalentes~:
  \begin{enumerate}[label=(\roman*)]
    \item $\cB$ est une base de $E$.
    \item Tout vecteur $v\in E$ s'écrit de manière \emph{unique} comme
      combinaison linéaire de $v_1,\ldots,v_n$~:
      \[
        v = \alpha_1 v_1 + \cdots + \alpha_n v_n,\qquad
        \alpha_1,\ldots,\alpha_n\in\K\text{ uniques.}
      \]
  \end{enumerate}
\end{theorem}

\begin{proof}
  \emph{(i)$\Rightarrow$(ii).} La famille est génératrice~: tout
  $v\in E$ s'écrit $v=\sum_i\alpha_i v_i$. Supposons deux
  décompositions~: $\sum_i\alpha_i v_i=\sum_i\beta_i v_i$. Alors
  $\sum_i(\alpha_i-\beta_i)v_i=\vzero$, et la liberté de $\cB$
  impose $\alpha_i-\beta_i=0$ pour tout $i$, d'où l'unicité.

  \emph{(ii)$\Rightarrow$(i).} L'existence de la décomposition implique
  que $\cB$ est génératrice. L'unicité implique la liberté~: si
  $\sum_i\alpha_i v_i=\vzero$, c'est la même décomposition que
  $\sum_i 0\cdot v_i=\vzero$, donc $\alpha_i=0$ pour tout $i$.
\end{proof}

\begin{definition}{Coordonnées}{def:coordonnees}
  \index{coordonnées}%
  Si $\cB=(v_1,\ldots,v_n)$ est une base de $E$ et si
  $v=\alpha_1 v_1+\cdots+\alpha_n v_n$, les scalaires
  $(\alpha_1,\ldots,\alpha_n)\in\K^n$ sont les \emph{coordonnées}
  (ou \emph{composantes}) de $v$ dans la base $\cB$.
\end{definition}

\begin{example}{Base canonique de $\K^n$}{ex:base-canonique}
  \index{base canonique}%
  La \emph{base canonique} de $\K^n$ est $\cE=(\vece{1},\ldots,\vece{n})$
  où $\vece{i}=(0,\ldots,0,1,0,\ldots,0)$ (le $1$ est en position $i$).
  Les coordonnées de $v=(x_1,\ldots,x_n)$ dans $\cE$ sont précisément
  $(x_1,\ldots,x_n)$.
\end{example}

\begin{example}{Base canonique de $\K_n[X]$}{ex:base-polynomes}
  La base canonique de $\K_n[X]$ est $(1,X,X^2,\ldots,X^n)$. Les
  coordonnées du polynôme $P=a_0+a_1 X+\cdots+a_n X^n$ sont
  $(a_0,a_1,\ldots,a_n)$.
\end{example}

\begin{example}{Base de $\Mat_{2}(\R)$}{ex:base-matrices}
  L'espace $\Mat_2(\R)$ des matrices $2\times 2$ admet pour base
  canonique~:
  \[
    E_{11}=\begin{pmatrix}1&0\\0&0\end{pmatrix},\;
    E_{12}=\begin{pmatrix}0&1\\0&0\end{pmatrix},\;
    E_{21}=\begin{pmatrix}0&0\\1&0\end{pmatrix},\;
    E_{22}=\begin{pmatrix}0&0\\0&1\end{pmatrix}.
  \]
  Les coordonnées de
  $A=\bigl(\begin{smallmatrix}a&b\\c&d\end{smallmatrix}\bigr)$ sont
  $(a,b,c,d)$.
\end{example}

% ============================================================
\section{Dimension}\label{sec:dimension}
% ============================================================

Le résultat fondamental suivant garantit que toutes les bases d'un
espace vectoriel ont le même nombre d'éléments.

\begin{lemma}{Lemme de Steinitz (lemme d'échange)}{lem:steinitz}
  \index{lemme de Steinitz}\index{lemme d'échange}%
  Soit $E$ un $\K$-espace vectoriel engendré par une famille
  $(g_1,\ldots,g_n)$ de $n$ vecteurs. Soit $(v_1,\ldots,v_p)$ une
  famille libre de $E$. Alors~:
  \begin{enumerate}[label=(\roman*)]
    \item $p\leq n$.
    \item On peut remplacer $p$ des $g_j$ par les $v_i$ pour obtenir
      une famille génératrice de $E$.
  \end{enumerate}
\end{lemma}

\begin{proof}
  On procède par récurrence sur $p$.

  \emph{Initialisation} ($p=0$)~: il n'y a rien à montrer.

  \emph{Hérédité.} Supposons le résultat vrai pour $p-1$ vecteurs
  libres, et soit $(v_1,\ldots,v_p)$ une famille libre de $E$. Par
  hypothèse de récurrence appliquée à $(v_1,\ldots,v_{p-1})$, on peut
  supposer (quitte à renuméroter) que
  $(v_1,\ldots,v_{p-1},g_p,\ldots,g_n)$ est génératrice de $E$, et
  $p-1\leq n$.

  Puisque cette famille engendre $E$, il existe
  $\alpha_1,\ldots,\alpha_{p-1},\beta_p,\ldots,\beta_n\in\K$ tels que
  \[
    v_p = \alpha_1 v_1+\cdots+\alpha_{p-1}v_{p-1}+\beta_p g_p
    +\cdots+\beta_n g_n.
  \]
  Affirmons qu'au moins un $\beta_j$ est non nul. En effet, si tous
  les $\beta_j$ étaient nuls, on aurait
  $v_p = \alpha_1 v_1+\cdots+\alpha_{p-1}v_{p-1}$, ce qui
  contredirait la liberté de $(v_1,\ldots,v_p)$.

  Soit $j_0\in\set{p,\ldots,n}$ tel que $\beta_{j_0}\neq 0$. On peut
  exprimer $g_{j_0}$ en fonction de $v_1,\ldots,v_p$ et des autres
  $g_j$ ($j\neq j_0$). En renumérotant, on suppose $j_0=p$, et on
  obtient que $(v_1,\ldots,v_p,g_{p+1},\ldots,g_n)$ engendre $E$.
  De plus, $p\leq n$ (car on a besoin d'au moins un $\beta_j$ non nul,
  ce qui requiert $n\geq p$).
\end{proof}

\begin{theorem}{Invariance de la dimension}{thm:invariance-dim}
  \index{dimension!invariance}%
  Si un $\K$-espace vectoriel $E$ admet une base finie, alors toutes
  ses bases ont le même nombre d'éléments.
\end{theorem}

\begin{proof}
  Soient $\cB=(v_1,\ldots,v_n)$ et $\cB'=(w_1,\ldots,w_m)$ deux bases
  de $E$. La famille $\cB'$ est libre et $\cB$ est génératrice, donc
  $m\leq n$ par le 
ef{lem:lem:steinitz}. Inversement, $\cB$ est libre
  et $\cB'$ est génératrice, donc $n\leq m$. D'où $m=n$.
\end{proof}

\begin{definition}{Dimension}{def:dimension}
  \index{dimension}%
  Un $\K$-espace vectoriel $E$ est dit \emph{de dimension finie}%
  \index{dimension finie} s'il admet une famille génératrice finie. Le
  nombre d'éléments de toute base de $E$ s'appelle la \emph{dimension}
  de $E$ et se note $\dim_\K E$ ou simplement $\dim E$.

  Par convention, $\dim\set{\vzero}=0$.
\end{definition}

\begin{theorem}{Existence de bases}{thm:existence-bases}
  Tout espace vectoriel de dimension finie admet une base. Plus
  précisément~:
  \begin{enumerate}[label=(\roman*)]
    \item Toute famille génératrice finie contient une base
      (on peut en extraire une sous-famille qui est une base).
    \item Toute famille libre peut être complétée en une base.
  \end{enumerate}
\end{theorem}

\begin{proof}
  \emph{(i)} Soit $(g_1,\ldots,g_n)$ une famille génératrice de $E$.
  Si elle est libre, c'est une base. Sinon, l'un des $g_j$ est
  combinaison linéaire des autres (par le 
ef{thm:thm:caract-liee})~; on
  peut le retirer et la famille restante engendre encore $E$. On
  itère ce procédé (qui termine car la famille diminue strictement à
  chaque étape) jusqu'à obtenir une famille libre, qui est alors une
  base.

  \emph{(ii)} Soit $(v_1,\ldots,v_p)$ une famille libre de $E$, et
  $(g_1,\ldots,g_n)$ une famille génératrice de $E$ (qui existe car
  $E$ est de dimension finie). On peut, par le procédé d'échange du
  
ef{lem:lem:steinitz}, compléter $(v_1,\ldots,v_p)$ en une famille
  génératrice en lui adjoignant certains des $g_j$. La famille obtenue
  est génératrice~; on en extrait une base par (i) en ne retirant
  aucun des $v_i$ (car ils forment une famille libre).
\end{proof}

\begin{example}{Dimensions classiques}{ex:dimensions}
  \begin{enumerate}[label=(\alph*)]
    \item $\dim\K^n = n$ (base canonique $(\vece{1},\ldots,\vece{n})$).
    \item $\dim\K_n[X] = n+1$ (base $(1,X,\ldots,X^n)$).
    \item $\dim\Mat_{n,p}(\K) = np$ (base $(E_{ij})_{1\leq i\leq n,\,1\leq j\leq p}$).
    \item $\K[X]$ est de dimension \emph{infinie} (la famille
      $(1,X,X^2,\ldots)$ est libre, donc $\dim\K[X]\geq n$ pour tout
      $n$).
  \end{enumerate}
\end{example}

\begin{theorem}{Théorème de la dimension pour les sous-espaces}{thm:dim-sev}
  Soit $E$ un $\K$-espace vectoriel de dimension finie $n$, et soit $F$
  un sous-espace vectoriel de $E$. Alors~:
  \begin{enumerate}[label=(\roman*)]
    \item $F$ est de dimension finie et $\dim F\leq\dim E$.
    \item $\dim F = \dim E$ si et seulement si $F=E$.
  \end{enumerate}
\end{theorem}

\begin{proof}
  \emph{(i)} Toute famille libre de $F$ est aussi une famille libre de
  $E$, donc a au plus $n=\dim E$ éléments (par le
  
ef{lem:lem:steinitz}). Soit $(v_1,\ldots,v_p)$ une famille libre de
  $F$ de cardinal maximal (un tel maximum existe car $p\leq n$).
  Montrons qu'elle engendre $F$~: si un vecteur $w\in F$ n'était pas
  dans $\Vect(v_1,\ldots,v_p)$, alors $(v_1,\ldots,v_p,w)$ serait
  libre dans $F$ (car $w\notin\Vect(v_1,\ldots,v_p)$), ce qui
  contredirait la maximalité de $p$. Donc
  $\Vect(v_1,\ldots,v_p)=F$, et $\cB$ est une base de $F$ avec
  $\dim F=p\leq n$.

  \emph{(ii)} Si $\dim F=\dim E=n$, soit $(v_1,\ldots,v_n)$ une base
  de $F$. C'est une famille libre de $n$ vecteurs dans $E$ de dimension
  $n$, donc c'est une base de $E$ (par la 
ef{prop:prop:libre-n-base}
  ci-dessous). D'où $F=\Vect(v_1,\ldots,v_n)=E$.
\end{proof}

\begin{proposition}{Critères de base en dimension $n$}{prop:libre-n-base}
  \index{critère de base}%
  Soit $E$ un $\K$-espace vectoriel de dimension $n$. Soit
  $(v_1,\ldots,v_n)$ une famille de $n$ vecteurs de $E$. Alors les
  assertions suivantes sont équivalentes~:
  \begin{enumerate}[label=(\roman*)]
    \item $(v_1,\ldots,v_n)$ est une base de $E$.
    \item $(v_1,\ldots,v_n)$ est libre.
    \item $(v_1,\ldots,v_n)$ est génératrice.
  \end{enumerate}
\end{proposition}

\begin{proof}
  \emph{(i)$\Rightarrow$(ii)} et \emph{(i)$\Rightarrow$(iii)} sont
  évidents.

  \emph{(ii)$\Rightarrow$(i).} Soit $\cB'$ une base de $E$ (de
  cardinal $n$). La famille $(v_1,\ldots,v_n)$ est libre et $\cB'$
  est génératrice de cardinal $n$. Par le 
ef{lem:lem:steinitz}, on
  peut remplacer les $n$ vecteurs de $\cB'$ par les $v_i$ (puisque
  $p=n$) pour obtenir une famille génératrice. Mais cette famille
  est précisément $(v_1,\ldots,v_n)$, qui est donc à la fois libre
  et génératrice~: c'est une base.

  \emph{(iii)$\Rightarrow$(i).} Si $(v_1,\ldots,v_n)$ est
  génératrice, on peut en extraire une base par le
  
ef{thm:thm:existence-bases}~(i). Cette base a $\dim E=n$ éléments.
  Puisqu'on part de $n$ vecteurs et qu'on ne peut en retirer aucun
  (la base extraite a aussi $n$ éléments), la famille de départ est
  elle-même une base.
\end{proof}

% ============================================================
\section{Dimension de la somme --- formule de Grassmann}\label{sec:grassmann}
% ============================================================

\begin{theorem}{Formule de Grassmann}{thm:grassmann}
  \index{formule de Grassmann}\index{Grassmann, formule de}%
  Soient $F$ et $G$ deux sous-espaces vectoriels de dimension finie
  d'un $\K$-espace vectoriel $E$. Alors $F+G$ est de dimension finie et
  \[
    \dim(F+G) = \dim F + \dim G - \dim(F\cap G).
  \]
\end{theorem}

\begin{proof}
  Posons $r=\dim(F\cap G)$, $p=\dim F$ et $q=\dim G$. Soit
  $(w_1,\ldots,w_r)$ une base de $F\cap G$.

  Puisque $F\cap G$ est un s.e.v.\ de $F$, on peut compléter cette
  famille en une base de $F$~: soit $(w_1,\ldots,w_r,u_1,\ldots,u_s)$
  une base de $F$, avec $s=p-r$.

  De même, soit $(w_1,\ldots,w_r,v_1,\ldots,v_t)$ une base de $G$,
  avec $t=q-r$.

  \emph{Affirmons que $\cB=(w_1,\ldots,w_r,u_1,\ldots,u_s,v_1,\ldots,v_t)$
  est une base de $F+G$.}

  \medskip
  \emph{$\cB$ est génératrice de $F+G$.} Tout $x\in F+G$ s'écrit
  $x=f+g$ avec $f\in F$ et $g\in G$. Or $f$ est combinaison linéaire
  des $w_i$ et des $u_j$, et $g$ est combinaison linéaire des $w_i$
  et des $v_k$. Donc $x$ est combinaison linéaire des $w_i$, $u_j$,
  $v_k$.

  \medskip
  \emph{$\cB$ est libre.} Supposons
  \[
    \sum_{i=1}^{r}\gamma_i w_i + \sum_{j=1}^{s}\alpha_j u_j
    + \sum_{k=1}^{t}\beta_k v_k = \vzero.
  \]
  Posons $h = \sum_{k=1}^{t}\beta_k v_k =
  -\sum_{i=1}^{r}\gamma_i w_i - \sum_{j=1}^{s}\alpha_j u_j$.
  Le membre droit appartient à $F$ (car les $w_i$ et $u_j$ sont dans
  $F$). Le membre gauche appartient à $G$ (car les $v_k$ sont dans
  $G$). Donc $h\in F\cap G$.

  Puisque $(w_1,\ldots,w_r)$ est une base de $F\cap G$, il existe
  $\delta_1,\ldots,\delta_r\in\K$ tels que
  $h=\sum_{i=1}^{r}\delta_i w_i$. D'où
  \[
    \sum_{k=1}^{t}\beta_k v_k - \sum_{i=1}^{r}\delta_i w_i = \vzero.
  \]
  Puisque $(w_1,\ldots,w_r,v_1,\ldots,v_t)$ est une base de $G$ (donc
  libre), on en déduit $\beta_1=\cdots=\beta_t=0$ et
  $\delta_1=\cdots=\delta_r=0$.

  Il reste $\sum_{i=1}^{r}\gamma_i w_i + \sum_{j=1}^{s}\alpha_j u_j
  = \vzero$. Puisque $(w_1,\ldots,w_r,u_1,\ldots,u_s)$ est une base
  de $F$ (donc libre), on a $\gamma_1=\cdots=\gamma_r=0$ et
  $\alpha_1=\cdots=\alpha_s=0$.

  Tous les coefficients sont nuls~: $\cB$ est libre.

  \medskip
  Ainsi $\cB$ est une base de $F+G$, et
  \[
    \dim(F+G) = r+s+t = r+(p-r)+(q-r) = p+q-r = \dim F+\dim G-\dim(F\cap G).
    \qedhere
  \]
\end{proof}

\begin{corollary}{Dimension de la somme directe}{cor:dim-somme-directe}
  Si $F\cap G=\set{\vzero}$ (somme directe), alors
  $\dim(F\oplus G)=\dim F+\dim G$.
\end{corollary}

\begin{proof}
  C'est le cas $\dim(F\cap G)=0$ de la formule de Grassmann.
\end{proof}

\begin{corollary}{Supplémentaire et dimension}{cor:supplementaire-dim}
  Soient $E$ un espace de dimension finie $n$ et $F$ un sous-espace de
  $E$ de dimension $p$. Tout supplémentaire $G$ de $F$ a pour dimension
  $n-p$.
\end{corollary}

\begin{proof}
  Si $E=F\oplus G$, alors $n=\dim E=\dim F+\dim G=p+\dim G$, d'où
  $\dim G=n-p$.
\end{proof}

\begin{corollary}{Condition de somme directe en dimension finie}{cor:condition-somme-directe}
  Soient $F$ et $G$ deux sous-espaces d'un espace $E$ de dimension
  finie. Alors $E=F\oplus G$ si et seulement si l'une des conditions
  suivantes est satisfaite~:
  \begin{enumerate}[label=(\roman*)]
    \item $E=F+G$ et $\dim E=\dim F+\dim G$~;
    \item $F\cap G=\set{\vzero}$ et $\dim E=\dim F+\dim G$.
  \end{enumerate}
\end{corollary}

\begin{proof}
  Dans les deux cas, la formule de Grassmann donne la condition manquante.

  \emph{(i)} Si $E=F+G$ et $\dim E=\dim F+\dim G$, la formule de
  Grassmann donne $\dim(F\cap G)=\dim F+\dim G-\dim(F+G)=\dim E-\dim E=0$,
  donc $F\cap G=\set{\vzero}$.

  \emph{(ii)} Si $F\cap G=\set{\vzero}$ et $\dim F+\dim G=\dim E$,
  alors $\dim(F+G)=\dim F+\dim G=\dim E$ par Grassmann, donc $F+G=E$
  par le 
ef{thm:thm:dim-sev}~(ii).
\end{proof}

% ============================================================
\section{Rang d'une famille de vecteurs}\label{sec:rang}
% ============================================================

\begin{definition}{Rang d'une famille de vecteurs}{def:rang-famille}
  \index{rang!d'une famille de vecteurs}%
  Soit $(v_1,\ldots,v_p)$ une famille de vecteurs d'un $\K$-espace
  vectoriel $E$ de dimension finie. Le \emph{rang} de cette famille est
  la dimension du sous-espace engendré~:
  \[
    \rg(v_1,\ldots,v_p) \deff \dim\Vect(v_1,\ldots,v_p).
  \]
\end{definition}

\begin{proposition}{Propriétés du rang}{prop:rang-proprietes}
  Soit $(v_1,\ldots,v_p)$ une famille de vecteurs de $E$.
  \begin{enumerate}[label=(\roman*)]
    \item $0\leq\rg(v_1,\ldots,v_p)\leq\min(p,\dim E)$.
    \item $\rg(v_1,\ldots,v_p)=p$ si et seulement si la famille est libre.
    \item $\rg(v_1,\ldots,v_p)=\dim E$ si et seulement si la famille est
      génératrice.
    \item $\rg(v_1,\ldots,v_p)=p=\dim E$ si et seulement si la famille
      est une base.
  \end{enumerate}
\end{proposition}

\begin{proof}
  \emph{(i)} Le rang est la dimension de $\Vect(v_1,\ldots,v_p)$, qui
  est un s.e.v.\ de $E$. Donc $\rg\leq\dim E$. De plus,
  $(v_1,\ldots,v_p)$ est une famille génératrice de
  $\Vect(v_1,\ldots,v_p)$, donc toute base extraite a au plus $p$
  éléments.

  \emph{(ii)} La famille est libre ssi elle est une base de
  $\Vect(v_1,\ldots,v_p)$, ce qui arrive ssi
  $\dim\Vect(v_1,\ldots,v_p)=p$.

  \emph{(iii)} La famille est génératrice ssi
  $\Vect(v_1,\ldots,v_p)=E$, soit $\rg=\dim E$.

  \emph{(iv)} Combine (ii) et (iii).
\end{proof}

% ============================================================
\section{Applications}\label{sec:ev-applications}
% ============================================================

\subsection{Espaces de solutions}\label{subsec:espaces-solutions}

\begin{proposition}{Dimension de l'espace des solutions}{prop:dim-solutions}
  \index{espace des solutions}%
  Soit $A\in\Mat_{n,p}(\K)$ une matrice de rang $r$. L'ensemble des
  solutions du système homogène $AX=0$ est un sous-espace vectoriel de
  $\K^p$ de dimension $p-r$.
\end{proposition}

Ce résultat fondamental, appelé \emph{théorème du rang}%
\index{théorème du rang}, sera démontré au 
ef{ch:espaces-vectoriels}
dans le cadre général des applications linéaires. Pour l'instant,
admettons-le et voyons-en une application.

\begin{example}{Espace des solutions d'un système}{ex:solutions-systeme}
  Considérons le système homogène sur $\R$~:
  \[
    \begin{cases}
      x+2y-z=0,\\
      2x+4y-2z=0.
    \end{cases}
  \]
  La matrice du système est
  $A=\bigl(\begin{smallmatrix}1&2&-1\\2&4&-2\end{smallmatrix}\bigr)$,
  de rang $1$ (la seconde ligne est le double de la première). L'espace
  des solutions est donc de dimension $3-1=2$. En résolvant, on trouve
  $x=-2y+z$, d'où~:
  \[
    \cS = \setcond{(-2y+z,\,y,\,z)}{y,z\in\R}
    = \Vect((-2,1,0),\,(1,0,1)).
  \]
  La famille $((-2,1,0),(1,0,1))$ est une base de $\cS$.
\end{example}

\subsection{Espaces de polynômes}\label{subsec:ev-polynomes}

\begin{example}{Polynômes pairs et impairs}{ex:polynomes-pairs-impairs}
  \index{polynôme pair}\index{polynôme impair}%
  Soit $E=\R_n[X]$ ($n\in\N$). Posons
  \begin{align*}
    F &= \setcond{P\in E}{P(-X)=P(X)} &&\text{(polynômes pairs)}, \\
    G &= \setcond{P\in E}{P(-X)=-P(X)} &&\text{(polynômes impairs)}.
  \end{align*}
  Alors $F$ et $G$ sont des sous-espaces de $E$, et $E=F\oplus G$.

  En effet, $F\cap G=\set{0}$ (si $P$ est à la fois pair et impair,
  $P(X)=-P(X)$ pour tout $X$, donc $2P=0$, d'où $P=0$ car
  $\car(\R)=0$). De plus, tout $P\in E$ s'écrit
  \[
    P(X) = \underbrace{\frac{P(X)+P(-X)}{2}}_{\in F}
    + \underbrace{\frac{P(X)-P(-X)}{2}}_{\in G}.
  \]
  Si $n=2m$, on a $\dim F=m+1$ et $\dim G=m$, d'où
  $\dim F+\dim G=(m+1)+m=2m+1=n+1=\dim E$.
  Si $n=2m+1$, on a $\dim F=m+1$ et $\dim G=m+1$, d'où
  $\dim F+\dim G=2m+2=n+1=\dim E$.
\end{example}

\subsection{Somme directe --- illustration géométrique}\label{subsec:somme-directe-geom}

\begin{center}
\begin{tikzpicture}[scale=1.3]
  % Subspace F (horizontal line)
  \draw[thick,blue!70!black,-{Stealth}] (-2.5,0) -- (2.5,0);
  \node[blue!70!black,below] at (2.3,0) {$F$};
  % Subspace G (line y=x)
  \draw[thick,red!70!black,-{Stealth}] (-1.5,-1.5) -- (1.8,1.8);
  \node[red!70!black,above left] at (1.7,1.7) {$G$};
  % A vector decomposition
  \coordinate (v) at (1.8,1);
  \coordinate (vF) at (1.8,0);
  \coordinate (vG) at (0.5,0.5);
  % Projection lines
  \draw[dashed,gray] (v) -- (vF);
  \draw[dashed,gray] (v) -- (0.5,0.5);
  \draw[dashed,gray] (0.5,0.5) -- (1.8,1);
  % Vectors
  \draw[-{Stealth},thick,purple!70!black] (0,0) -- (v)
    node[above right,font=\small] {$w=u+v$};
  \draw[-{Stealth},thick,blue!70!black] (0,0) -- (vF)
    node[below,font=\small] {$u\in F$};
  \draw[-{Stealth},thick,red!70!black] (0,0) -- (vG)
    node[left,font=\small] {$v\in G$};
  % Parallelogram
  \draw[dotted,gray] (vF) -- (v);
  \draw[dotted,gray] (vG) -- (v);
  % Origin
  \fill (0,0) circle (1.5pt) node[below left] {$\vzero$};
  % Label
  \node[font=\small] at (0,-2) {$\R^2 = F\oplus G$};
\end{tikzpicture}
\captionof{figure}{Décomposition d'un vecteur selon deux sous-espaces
  supplémentaires dans $\R^2$.}
\label{fig:somme-directe-R2}
\end{center}

% ============================================================
\section{Exercices}\label{sec:exo-ch1}
% ============================================================

\begin{exercise}{\basic{} Vérification des axiomes}{exo:verif-axiomes}
  Montrer que l'ensemble $E=\R^2$ muni des opérations
  \[
    (x_1,y_1)\oplus(x_2,y_2)=(x_1+x_2,\,y_1+y_2),\qquad
    \alpha\odot(x,y)=(\alpha x,\,\alpha y)
  \]
  est un $\R$-espace vectoriel. Identifier le vecteur nul et l'opposé
  de $(x,y)$.
\end{exercise}

\begin{exercise}{\basic{} Espace vectoriel ou non~?}{exo:ev-ou-non}
  Pour chacun des ensembles suivants, munis des opérations usuelles
  d'addition et de multiplication scalaire, dire s'il s'agit d'un
  $\R$-espace vectoriel. Justifier.
  \begin{enumerate}[label=(\alph*)]
    \item $E_1=\setcond{(x,y)\in\R^2}{x+y=0}$.
    \item $E_2=\setcond{(x,y)\in\R^2}{x+y=1}$.
    \item $E_3=\setcond{(x,y)\in\R^2}{xy=0}$.
    \item $E_4=\setcond{P\in\R[X]}{P(0)=0}$.
    \item $E_5=\setcond{f\in\cC^0(\R,\R)}{f(0)=0\text{ et }f(1)=0}$.
  \end{enumerate}
\end{exercise}

\begin{exercise}{\basic{} Sous-espace vectoriel}{exo:sev-basique}
  Montrer que les ensembles suivants sont des sous-espaces vectoriels
  de l'espace ambiant indiqué~:
  \begin{enumerate}[label=(\alph*)]
    \item $F=\setcond{(x,y,z)\in\R^3}{x-2y+z=0}$ dans $\R^3$.
    \item $G=\setcond{A\in\Mat_2(\R)}{\tr(A)=0}$ dans $\Mat_2(\R)$.
    \item $H=\setcond{P\in\R_3[X]}{P(1)=0}$ dans $\R_3[X]$.
  \end{enumerate}
\end{exercise}

\begin{exercise}{\basic{} Famille libre ou liée}{exo:libre-liee}
  Déterminer si les familles suivantes sont libres ou liées. Si la
  famille est liée, exhiber une relation de dépendance linéaire.
  \begin{enumerate}[label=(\alph*)]
    \item $((1,2),(3,6))$ dans $\R^2$.
    \item $((1,1,0),(0,1,1),(1,0,-1))$ dans $\R^3$.
    \item $(1,1+X,1+X+X^2)$ dans $\R_2[X]$.
    \item $(e^x,e^{2x},e^{3x})$ dans $\cC^0(\R,\R)$.
  \end{enumerate}
\end{exercise}

\begin{exercise}{\basic{} Base et coordonnées}{exo:base-coord}
  \begin{enumerate}[label=(\alph*)]
    \item Vérifier que $\cB=((1,1),(1,-1))$ est une base de $\R^2$.
    \item Exprimer le vecteur $v=(3,5)$ dans la base $\cB$ (trouver ses
      coordonnées).
    \item Montrer que $\cB'=((1,1,1),(1,1,0),(1,0,0))$ est une base de
      $\R^3$ et trouver les coordonnées de $(2,3,5)$ dans $\cB'$.
  \end{enumerate}
\end{exercise}

\begin{exercise}{\intermediate{} Intersection et somme}{exo:inter-somme}
  Dans $\R^4$, soient
  \begin{align*}
    F &= \Vect((1,0,1,0),\,(0,1,0,1)), \\
    G &= \Vect((1,1,0,0),\,(0,0,1,1)).
  \end{align*}
  \begin{enumerate}[label=(\alph*)]
    \item Déterminer $\dim F$ et $\dim G$.
    \item Déterminer une base de $F\cap G$.
    \item En déduire $\dim(F+G)$ par la formule de Grassmann.
    \item La somme $F+G$ est-elle directe~?
  \end{enumerate}
\end{exercise}

\begin{exercise}{\intermediate{} Somme directe de sous-espaces}{exo:somme-directe}
  Dans $\R^3$, soient $F=\setcond{(x,y,z)}{x+y=0}$ et
  $G=\Vect((1,1,1))$.
  \begin{enumerate}[label=(\alph*)]
    \item Montrer que $F$ est un sous-espace de dimension $2$ et
      déterminer une base de $F$.
    \item Montrer que $F\cap G=\set{\vzero}$.
    \item En déduire que $\R^3=F\oplus G$.
    \item Décomposer le vecteur $(2,3,5)\in\R^3$ en $u+v$ avec $u\in F$
      et $v\in G$.
  \end{enumerate}
\end{exercise}

\begin{exercise}{\intermediate{} Dimension d'un espace de matrices}{exo:dim-matrices}
  Soit $E=\Mat_n(\R)$, et soient
  \begin{align*}
    \mathcal{S}_n(\R) &= \setcond{A\in E}{\transpose{A}=A}
    &&\text{(matrices symétriques)},\\
    \mathcal{A}_n(\R) &= \setcond{A\in E}{\transpose{A}=-A}
    &&\text{(matrices antisymétriques)}.
  \end{align*}
  \begin{enumerate}[label=(\alph*)]
    \item Montrer que $\mathcal{S}_n(\R)$ et $\mathcal{A}_n(\R)$ sont
      des sous-espaces de $E$.
    \item Calculer $\dim\mathcal{S}_n(\R)$ et $\dim\mathcal{A}_n(\R)$.
    \item Montrer que $E=\mathcal{S}_n(\R)\oplus\mathcal{A}_n(\R)$.
  \end{enumerate}
\end{exercise}

\begin{exercise}{\intermediate{} Sous-espace engendré et rang}{exo:rang}
  On considère les vecteurs de $\R^4$~:
  \[
    v_1=(1,2,1,0),\quad v_2=(1,1,0,1),\quad
    v_3=(3,5,2,1),\quad v_4=(0,1,1,-1).
  \]
  \begin{enumerate}[label=(\alph*)]
    \item Montrer que $(v_1,v_2,v_3,v_4)$ est liée en exhibant une
      relation de dépendance.
    \item Déterminer le rang de cette famille.
    \item En extraire une base de $\Vect(v_1,v_2,v_3,v_4)$.
  \end{enumerate}
\end{exercise}

\begin{exercise}{\intermediate{} Formule de Grassmann}{exo:grassmann}
  Dans $\R^4$, soient
  \begin{align*}
    F &= \setcond{(x,y,z,t)\in\R^4}{x+y=0,\;z-t=0}, \\
    G &= \setcond{(x,y,z,t)\in\R^4}{x-z=0,\;y+t=0}.
  \end{align*}
  \begin{enumerate}[label=(\alph*)]
    \item Déterminer des bases de $F$ et $G$, et calculer leurs
      dimensions.
    \item Déterminer une base de $F\cap G$ et calculer sa dimension.
    \item Vérifier la formule de Grassmann.
    \item La somme $F+G$ est-elle directe~? Est-elle égale à $\R^4$~?
  \end{enumerate}
\end{exercise}

\begin{exercise}{\challenging{} Supplémentaire dans $\R_n[X]$}{exo:supplementaire-polynomes}
  Soient $n\geq 1$ et $F=\setcond{P\in\R_n[X]}{P(0)=P(1)=0}$.
  \begin{enumerate}[label=(\alph*)]
    \item Montrer que $F$ est un sous-espace de $\R_n[X]$ et calculer
      sa dimension.
    \item Trouver un supplémentaire $G$ de $F$ dans $\R_n[X]$ et en
      donner une base.
    \item Pour $n=3$, décomposer le polynôme $P=1+X+X^2+X^3$ en $f+g$
      avec $f\in F$ et $g\in G$.
  \end{enumerate}
\end{exercise}

\begin{exercise}{\challenging{} Dimension infinie}{exo:dim-infinie}
  \begin{enumerate}[label=(\alph*)]
    \item Montrer que $\K[X]$ est de dimension infinie.
    \item Montrer que l'espace $\cC^0([0,1],\R)$ des fonctions continues
      sur $[0,1]$ est de dimension infinie.

      \emph{Indication~:} montrer que les fonctions $1,x,x^2,\ldots$
      forment une famille libre.
    \item Montrer que l'espace $\R^\N$ des suites réelles est de
      dimension infinie.
  \end{enumerate}
\end{exercise}

\begin{exercise}{\challenging{} Une somme qui n'est pas directe}{exo:somme-non-directe}
  Dans $\R^3$, soient $F_1=\Vect((1,0,0),(0,1,0))$,
  $F_2=\Vect((0,1,0),(0,0,1))$ et $F_3=\Vect((1,0,0),(0,0,1))$.
  \begin{enumerate}[label=(\alph*)]
    \item Montrer que $F_1+F_2+F_3=\R^3$.
    \item Montrer que $F_i\cap F_j=\set{\vzero}$ n'est \emph{pas}
      vérifié pour tous $i\neq j$.
    \item La somme $F_1+F_2+F_3$ est-elle directe~? Justifier.
    \item Comparer avec la condition de la
      
ef{prop:prop:caract-somme-directe-n}.
  \end{enumerate}
\end{exercise}

\begin{exercise}{\challenging{} Lemme de Steinitz appliqué}{exo:steinitz-applique}
  Soit $E=\R^3$ et $\cB=(e_1,e_2,e_3)$ la base canonique. Soit
  $v_1=(1,1,0)$, $v_2=(0,1,1)$.
  \begin{enumerate}[label=(\alph*)]
    \item Vérifier que $(v_1,v_2)$ est libre.
    \item Appliquer le procédé du 
ef{lem:lem:steinitz} pour compléter
      $(v_1,v_2)$ en une base de $\R^3$ en utilisant les vecteurs de
      $\cB$.
    \item Combien de façons y a-t-il de choisir le troisième vecteur
      parmi $\set{e_1,e_2,e_3}$ pour obtenir une base~?
  \end{enumerate}
\end{exercise}

% ============================================================
\section*{Résumé du chapitre}
\addcontentsline{toc}{section}{Résumé du chapitre}
% ============================================================

\begin{framed}
\textbf{Ce qu'il faut retenir.}

\begin{enumerate}[label=\textbf{\arabic*.},leftmargin=2em]
  \item \textbf{Espace vectoriel.} Un $\K$-espace vectoriel est un
    ensemble $E$ muni d'une addition (qui en fait un groupe abélien) et
    d'une multiplication par un scalaire de $\K$ satisfaisant les
    axiomes de compatibilité, d'élément unité et de distributivité.

  \item \textbf{Sous-espace vectoriel.} $F\subseteq E$ est un s.e.v.\
    ssi $F\neq\emptyset$ et $F$ est stable par combinaison linéaire.

  \item \textbf{Intersection et somme.} L'intersection de sous-espaces
    est un s.e.v. La somme $F+G$ est le plus petit s.e.v.\ contenant
    $F$ et $G$. La réunion n'est en général \emph{pas} un s.e.v.

  \item \textbf{Somme directe.} $F+G$ est directe ssi
    $F\cap G=\set{\vzero}$. On note $F\oplus G$. Si $E=F\oplus G$,
    on dit que $F$ et $G$ sont supplémentaires.

  \item \textbf{Famille libre.} $(\alpha_1 v_1+\cdots+\alpha_n v_n
    =\vzero) \implies (\alpha_i=0\;\forall i)$.

  \item \textbf{Famille génératrice.} Tout vecteur de $E$ est
    combinaison linéaire de la famille.

  \item \textbf{Base.} Famille à la fois libre et génératrice. Tout
    vecteur s'y décompose de manière unique (coordonnées). En dimension
    finie $n$~: famille libre de $n$ vecteurs $\iff$ famille
    génératrice de $n$ vecteurs $\iff$ base.

  \item \textbf{Dimension.} Nombre d'éléments de toute base (invariant).
    $\dim\K^n=n$, $\dim\K_n[X]=n+1$, $\dim\Mat_{n,p}(\K)=np$.

  \item \textbf{Formule de Grassmann.}
    $\dim(F+G)=\dim F+\dim G-\dim(F\cap G)$.

  \item \textbf{Rang.} Le rang d'une famille $(v_1,\ldots,v_p)$ est
    $\dim\Vect(v_1,\ldots,v_p)$.

  \item \textbf{Théorèmes-clés.}
    \begin{itemize}
      \item Lemme de Steinitz (échange) --- une famille libre a au plus
        autant d'éléments qu'une famille génératrice.
      \item Toute famille libre peut être complétée en une base.
      \item $\dim F\leq\dim E$, avec égalité ssi $F=E$.
    \end{itemize}
\end{enumerate}
\end{framed}
